\section{Conclusion}

Managed freight lanes are the correct primary capacity investment for T1 corridors because they solve the structural failure mode of the US interstate system: the progressive collapse of freight performance as induced demand fills any new general purpose capacity within a decade. The throughput analysis in Section~\ref{sec:throughput} demonstrates that seven of the eight T1 corridors currently operate in chronic LOS F during peak periods, with PTI values of 1.58--2.21 that make shipper SLA commitments economically impossible. The I-40 exception --- operating at peak V/C 0.84 with PTI 1.11 --- demonstrates that the problem is not T1 classification per se but demand-capacity imbalance, and that the I2.0 investment program correctly directs managed lane resources only where the imbalance exists.

\subsection*{The Self-Reinforcing Benefit Structure}

The managed freight lane investment case rests on four compounding benefits that together create a return structure unlike any other T1 investment category:

\textbf{Dedicated capacity with permanent benefit.} Unlike general purpose lane additions, which induce demand and restore original V/C ratios within 10 years \citep{Duranton2011}, freight-only managed lanes with access control cannot be colonized by induced passenger demand. Induced demand fills the capacity that new vehicles can access; if the new lanes are inaccessible to passenger vehicles by design, passenger demand cannot consume them. The 45,600 commercial vehicle equivalent capacity addition per corridor per day remains a freight capacity addition in year 20 as in year 1. This is the core permanent infrastructure argument that justifies the \$177--265 billion program cost.

\textbf{Platooning enablement and fuel savings that fund operations.} Access-controlled freight lanes eliminate the passenger vehicle interactions --- lane changes, speed differentials, merge conflicts --- that prevent platooning in mixed traffic. At 65 mph sustained speed with consistent following distances, truck platooning achieves 20--35\% aerodynamic drag reduction for trailing vehicles \citep{Browand2004}. At current diesel prices and commercial vehicle mileage, this translates to approximately \$0.08 per mile in fuel savings per following vehicle. On I-75 Atlanta, with approximately 8,000 freight vehicles per day on the managed lanes, the annual fuel savings of \$180 million approaches the annualized construction cost of the corridor segment --- a self-funding mechanism not available to any other T1 investment category.

\textbf{PTI improvement enabling shipper SLA capability.} The reduction from current PTI of 1.58--2.21 to a managed-lane PTI target of $\leq$1.15 is not an operational refinement; it is a market-enabling change. At current PTI values, shippers cannot offer 95th-percentile on-time guarantees without padding transit windows by 58--121\% above free-flow. This padding absorbs the premium that time-sensitive freight should command, compressing shipper margins and suppressing modal share for truck freight on time-sensitive commodities. At PTI 1.15, a 15\% schedule buffer is sufficient for 95th-percentile delivery confidence --- a threshold at which pharmaceutical, perishable, and automotive JIT freight can be priced and contractually committed without the hedging discount \citep{ATRI_costs2024}.

\textbf{GP lane LOS improvement without new GP lane construction.} Transferring 22\% of AADT (the approximate commercial vehicle share on high-volume T1 segments) to dedicated freight lanes reduces GP lane effective V/C from breakdown-threshold levels to LOS D--E. This improvement occurs without adding a single general purpose lane --- which would itself induce demand, filling within a decade and delivering no permanent benefit. The GP lane LOS improvement is a political and fiscal externality of the managed lane investment: passenger vehicle users benefit, and no additional public capital is required to produce that benefit.

\subsection*{The NPV Case on the Highest-Volume Segments}

The freight economics justify the investment on the highest-volume T1 segments without relying on any platooning premium or GP lane spillover benefit. The computation proceeds from ATRI truck-operating cost data \citep{ATRI_costs2024}:

At \$150 per truck-hour ATRI operating cost and 45,600 commercial vehicle equivalents per day of added capacity, a conservative capacity utilization assumption of 70\% yields 31,920 truck-equivalents per day of realized freight benefit. Each truck-hour of delay eliminated on the managed lanes reduces cost by \$150. At the current peak-period delay differential between managed and GP lanes --- estimated at 0.8--1.2 hours per 100-mile segment at current LOS F conditions --- the daily freight benefit per 100-mile managed lane segment is 31,920 trucks $\times$ 1.0 hour $\times$ \$150 = \$4.8 million per day, or approximately \$1.75 billion per year per 100-mile segment.

Against an annualized construction cost of \$150--225 million per 100-mile segment (at \$10--15M per lane-mile, 2 lanes per direction, 30-year amortization at 4\% discount rate), the simple benefit-cost ratio at the highest-volume corridors is 8:1 to 12:1. The managed lane investment pays back the annualized construction cost within 8--12 years on the I-10, I-90, and I-80 segments that show peak V/C above 2.0. The NY--LA corridor alone --- where managed lanes reduce transit time by 22\% (one full calendar day at HOS-compliant driving schedules) --- generates an estimated \$10.5 billion in annual shipper benefit across the primary O-D pair \citep{FMCSA_HOS2024}.

\subsection*{The I-40 Exception Proves the Rule}

Interstate 40 at V/C 0.84 is the one T1 corridor correctly excluded from the managed freight lane program. Its exclusion demonstrates that the I2.0 methodology is demand-driven: tier classification identifies strategic corridors, but capacity analysis determines whether managed lanes are the right investment. A strategically critical corridor with adequate capacity does not need managed lanes --- it needs resilience investment (compound exposure remediation per Paper B.3 \citep{ROUTE_B3}) and access improvements (EV charging, modern rest areas, intermodal integration). The tier does not mandate the investment type; the performance data does.

This distinction has policy significance. The managed freight lane program is sometimes presented as a universal T1 upgrade --- applying the same treatment to every Primary Artery regardless of current performance. The I-40 exclusion, supported by convergent evidence from HPMS V/C analysis and ATRI revealed-preference data showing zero top-50 bottleneck locations on I-40 \citep{ROUTE_B2}, demonstrates that the program correctly discriminates between corridors that need capacity and corridors that need other interventions. This discrimination is a property of the underlying methodology, not a post-hoc adjustment: the tier classification and the capacity analysis are independent datasets that converge on the same conclusion.

\subsection*{Implementation Requirement: Access Control Permanence}

The permanent benefit of managed freight lanes depends entirely on the permanence of access control. A managed freight lane that becomes a general purpose lane --- whether through regulatory erosion, political pressure, or informal non-enforcement --- loses its capacity benefit through induced demand within the same 10-year window as any general purpose lane addition. Access control must be:

\begin{enumerate}
\item \textbf{Physical}: separated by concrete barrier, not paint; no passenger vehicle on-ramps in the managed lane footprint
\item \textbf{Enforced}: automatic license plate recognition and weigh-in-motion enforcement at all access points, with real-time violation detection
\item \textbf{Statutory}: federal designation as freight-only infrastructure, requiring affirmative federal action to modify --- not subject to state administrative reclassification
\end{enumerate}

Without all three conditions, the managed freight lane investment is a throughput loan, not a throughput asset. The \$177--265 billion program cost is justified only under the assumption that the capacity benefit is permanent. The architecture of access control --- physical, enforced, statutory --- is the single implementation requirement on which the entire NPV case depends.

\subsection*{Future Work}

Three extensions are warranted. First, the platooning benefit estimate uses aggregate aerodynamic drag data \citep{Browand2004}; fleet-composition-weighted platooning benefits by corridor (accounting for trailer type, load factor, and penetration of platooning-capable trucks in the commercial fleet) would improve the precision of the fuel savings estimate. Second, the PTI model uses historical incident data as a proxy for future managed-lane PTI; a simulation-based approach incorporating the incident-reduction effect of truck removal from GP lanes would strengthen the PTI projection. Third, the NY--LA transit time analysis uses a single representative route; O-D pair analysis across the full ROUTE freight demand matrix \citep{ROUTE_C1} would quantify the managed lane benefit across all T1 corridor O-D pairs, not just the primary transcontinental pair.
